\documentclass[11pt,a4paper]{article}

\usepackage{fontspec}
\setmainfont{Libertinus Serif}
\usepackage[greek.ancient,spanish]{babel}
\babelfont[greek]{rm}{Libertinus Serif}
\usepackage[a4paper,margin=3.2cm]{geometry}
\usepackage{microtype}
\usepackage[hidelinks]{hyperref}

% Griego politónico escrito directamente en UTF-8:
% \gr{...} para palabras sueltas; bloques con \foreignlanguage{greek}{...}
\newcommand{\gr}[1]{\foreignlanguage{greek}{#1}}
\newcommand{\stanzabreak}{\par\vspace{0.6\baselineskip}}

\setlength{\parindent}{1.25em}
\setlength{\parskip}{0.35em}
\raggedbottom


\title{Informe sobre las referencias a Nolan}
\author{}
\date{}

\newcommand{\justificada}{\textbf{Justificada.}}
\newcommand{\parcial}{\textbf{Parcialmente justificada.}}
\newcommand{\prescindible}{\textbf{Prescindible.}}

\begin{document}

\maketitle

\section*{Diagnóstico}

Nolan funciona bien como encarnación de «la \emph{Odisea} según Hollywood» cuando la comparación revela una pérdida de ambigüedad, crueldad, inteligencia o agencia respecto a Homero. El libro deja de sostener esa tesis cuando juzga el reparto, acumula omisiones sin consecuencias o lleva una cuenta de aciertos y errores cinematográficos. La solución no es retirar a Nolan, sino reducir los pasajes que personalizan en él un problema que el propio libro define como industrial y cultural.

No se cuentan las referencias actualmente comentadas, que no aparecen en el libro compilado.

\section*{Casos}

\begin{enumerate}

\item \textbf{«De libros y películas», líneas 43--57 y 75: cine, literatura y cultura audiovisual.} \justificada Es la premisa del libro y deja claro que Nolan es un caso de estudio, no el verdadero acusado. Mantener.

\item \textbf{«De libros y películas», líneas 59--63: histrionismo, estética \emph{new age}, buenismo y Curro Romero.} \parcial La expectativa ante una adaptación difícil es pertinente; el juicio sobre los tics personales del director acerca el pasaje a una crítica de cine. Conviene abreviarlo.

\item \textbf{«De libros y películas», líneas 69--73: críticas irrelevantes y Ulises con estrés postraumático.} \justificada El primer pasaje fija correctamente qué no interesa discutir; el segundo contrapone dos lecturas del héroe. Mantener.

\item \textbf{«El cine teme a Homero», líneas 49--51: Héctor en \emph{Troya} y la adaptación como traducción.} \justificada Extiende el argumento más allá de Nolan y demuestra que el problema pertenece al cine de masas. Mantener.

\item \textbf{«El cine teme a Homero», líneas 53--65: balance general de la \emph{Nodisea}.} \justificada Funciona como mapa de los ejemplos posteriores y reconoce los aciertos de la película. Mantener; no añadir más ejemplos aquí.

\item \textbf{«El cine teme a Homero», líneas 67--69: incoherencia de la \emph{xenía} y riesgo de que la película se convierta en canon.} \justificada Es una de las formulaciones centrales de la tesis. Mantener.

\item \textbf{«Una osada Odisea», línea 9: Nolan llena las librerías de clásicos.} \justificada No demuestra la tesis, pero explica el origen material del libro y reconoce un efecto beneficioso del estreno. Mantener.

\item \textbf{«Helena», líneas 56--58: supresión de los dioses.} \justificada Plantea una dificultad real del cine contemporáneo y la conecta con la reducción psicológica de lo sobrenatural. Mantener.

\item \textbf{«Helena», líneas 64, 143 y 306--307: de semidiosa impune a mujer maltratada.} \justificada Es un ejemplo nítido de sustitución de un personaje ambiguo por un cliché moral reconocible. Mantener.

\item \textbf{«Helena», líneas 182, 225 y 249: el mendigo, Nadie y el caballo.} \parcial La omisión de las astucias de Odiseo y la sustitución de la voz de Helena por acción sostienen la tesis. La larga nota sobre las ruedas del caballo repite una cuestión que la introducción ya declara irrelevante; conviene recortarla.

\item \textbf{«Menelao», línea 3: reinterpretar no equivale a simplificar.} \justificada Es una precisión necesaria: el libro no exige fidelidad literal, sino una reescritura con ambición. Mantener.

\item \textbf{«Calipso», líneas 47, 117 y 134: dificultad de representar a los dioses.} \justificada Generaliza el problema y admite por qué Hollywood prefiere eliminarlos. Mantener.

\item \textbf{«Calipso», líneas 70 y 239: Atenea en Troya y la soledad de Calipso.} \justificada Son contraejemplos breves que muestran que la película puede captar aspectos de Homero. La mención de Charlize Theron es secundaria, pero no llega a desviar el capítulo.

\item \textbf{«Calipso», línea 281: loto, culpa y estrés postraumático.} \justificada Analiza una reinterpretación coherente y señala con precisión qué cambia en el personaje. Mantener.

\item \textbf{«Nausícaa», línea 59: la polémica sobre \emph{dad}.} \parcial Sirve para rechazar una crítica trivial y conduce al \emph{páppa} homérico. Mantener, porque el centro del pasaje acaba siendo Homero y no Nolan.

\item \textbf{«Esqueria», líneas 52--54: eliminación de Nausícaa y los feacios.} \justificada Es probablemente el mejor ejemplo del libro: la película suprime el principal contraejemplo a su propia tesis sobre la \emph{xenía}. Mantener íntegro.

\item \textbf{«Bardos», líneas 3 y 67: Travis Scott como Demódoco y la inteligencia del espectador.} \parcial La reflexión final sobre la sutileza de Homero está plenamente justificada. La fantasía inicial sobre el reparto repite la ausencia de Esqueria y puede reducirse a una sola frase.

\item \textbf{«La Gran Guerra», línea 105: Odiseo convertido en conscripto moral.} \justificada Contrasta directamente la responsabilidad del héroe homérico con su exculpación hollywoodiense. Mantener.

\item \textbf{«Cíclope», líneas 3--7: balance jocoso de agravios y omisión de los lotófagos.} \prescindible Es un marcador de crítica cinematográfica y repite asuntos ya tratados. Conviene reducirlo al único dato necesario: Nolan elimina la estratagema de Nadie.

\item \textbf{«Cíclope», línea 256: Polifemo y la \emph{xenía}.} \justificada Une caracterización y estructura: al volver simple al cíclope, la película pierde el caso más consciente de violación de la hospitalidad. Mantener.

\item \textbf{«Cíclope», líneas 289, 355 y 531: peñasco, embriaguez y desenlace con Poseidón.} \justificada Los detalles no se usan para puntuar la fidelidad, sino para mostrar cómo se reduce la astucia y la arrogancia de Odiseo. Mantener; la breve aprobación del peñasco basta como contrapeso.

\item \textbf{«La bolsa de los vientos», líneas 80, 155 y 174: compañeros y lestrigones.} \justificada Estos aciertos prueban que el problema no es la incapacidad técnica de Nolan: la acción y la economía narrativa funcionan cuando no exigen ambigüedad moral. Mantener.

\item \textbf{«La isla de la hechicera», líneas 5--7, 43 y 78: Circe plana y transformación manual.} \justificada El contraste entre espectáculo plástico y pobreza del personaje sostiene la tesis. La cirugía de los cerdos se comenta dos veces; basta una.

\item \textbf{«Hades», líneas 3 y 60: excelencia visual y película de zombis.} \justificada La apertura es un contraejemplo importante. La comparación con los zombis es una digresión expresiva, pero vuelve a dirigir la atención hacia la potencia visual de Homero y puede mantenerse.

\item \textbf{«Hades», línea 153: profecía inevitable frente a prueba superable.} \justificada La diferencia afecta al libre albedrío y no solo a la fidelidad argumental. Mantener.

\item \textbf{«Hades», línea 242: Nolan se salta a Aquiles.} \prescindible La propia nota admite que es una exigencia de guion y no extrae ninguna consecuencia interpretativa. Suprimir esa observación.

\item \textbf{«Agamenón y Clitemnestra», líneas 3 y 134: retrato, atuendo y doble reparto.} \parcial Reconocer que la adaptación funciona es útil; el atuendo Darth Vader y el doble papel de Lupita Nyong'o pertenecen más a una reseña. Puede conservarse el juicio general y abreviar los detalles.

\item \textbf{«Agamenón y Clitemnestra», línea 209: justificación del sacrificio de Ifigenia.} \justificada Señala otra operación de moralización y exculpación psicológica. Mantener.

\item \textbf{«Sirenas», líneas 3--45.} \justificada Es la excepción deliberada: el diálogo con Nolan constituye la arquitectura del capítulo y culmina en un caso donde el cine supera momentáneamente a la palabra. Mantener íntegro.

\item \textbf{«El arco funesto», líneas 8--17: domesticación de Telémaco y Odiseo.} \justificada Presenta el núcleo del problema: Hollywood no solo reduce la violencia, sino que altera la naturaleza de ambos personajes. Mantener.

\item \textbf{«El arco funesto», líneas 34 y 212: escenas bien resueltas.} \parcial La primera aprobación equilibra el juicio antes de señalar la omisión de Telémaco. «Es otra escena muy bien conseguida» no añade nada y puede desaparecer.

\item \textbf{«El arco funesto», línea 185: Penélope sabe quién es el mendigo.} \justificada Es un ejemplo excelente de pérdida de agencia femenina por simplificación del argumento. Mantener.

\item \textbf{«El arco funesto», línea 227: el arco como advertencia moral.} \justificada La película convierte un gesto de dominio en una oportunidad misericordiosa y vuelve a santificar al héroe. Mantener.

\item \textbf{«La matanza de los pretendientes», líneas 3 y 71: Antínoo y la oferta de compensación.} \justificada La omisión del perdón rechazado demuestra con claridad qué violencia no tolera Hollywood. Mantener.

\item \textbf{«La matanza de los pretendientes», líneas 100, 155 y 183: armas, reparto y Melantio.} \parcial La reducción del peligro y el apartamiento de Telémaco son pertinentes. La comparación física entre Damon y Cavill es crítica de reparto y contradice el desinterés declarado por esas cuestiones; suprimirla.

\item \textbf{«La matanza de los pretendientes», línea 406: lobo con piel de oveja.} \justificada Resume el criterio del libro: toda reescritura es lícita, pero no toda simplificación resulta fértil. Mantener.

\item \textbf{«El asesinato de las niñas», líneas 34 y 206--210: lo que Hollywood no puede mostrar.} \justificada Es la demostración culminante de la tesis y desplaza expresamente la acusación desde Nolan hacia el cine de masas. Mantener íntegro.

\item \textbf{«El asesinato de las niñas», línea 204: \emph{Troya}, Príamo y Aquiles.} \justificada El contraste demuestra que Hollywood admite la violencia cuando puede resolverla mediante compasión y nobleza. Mantener.

\item \textbf{«Penélope», línea 108: reencuentro edulcorado.} \justificada La comparación descubre una pérdida de tiempo vivido, extrañeza y cautela, no una mera diferencia de puesta en escena. Mantener.

\item \textbf{«El último canto», línea 17: ninguna película muestra a las esclavas.} \justificada Amplía el caso de las muchachas desde Nolan al mecanismo general de ocultación de la violencia. Mantener.

\item \textbf{«El último canto», línea 398: navegación hacia el sol poniente.} \parcial El contraste entre el cierre grandilocuente y la brusquedad de Homero sirve a la tesis sobre la moralización. «Ridículas frases» vuelve el pasaje innecesariamente reseñístico; conviene conservar el contraste y rebajar el adjetivo.

\end{enumerate}

\section*{Prioridad editorial}

Las referencias decisivas son las relativas a Helena, la \emph{xenía}, Esqueria, Polifemo, Circe, Telémaco, Penélope y el asesinato de las muchachas. Los principales recortes deberían concentrarse en cuatro lugares: el balance inicial de «Cíclope», la omisión de Aquiles en «Hades», los comentarios de reparto en «Agamenón y Clitemnestra» y «La matanza de los pretendientes», y las aprobaciones genéricas repetidas en «El arco funesto». Con esos ajustes, Nolan seguiría presente en todo el recorrido, pero como síntoma de Hollywood y no como objeto de una reseña prolongada.

\end{document}
